Low-rank adaptation~\citep{hu2022lora} has made it cheap to produce many
specialised versions of one base model, and correspondingly attractive to
fold several of them back into a single set of weights. A large family of
merging heuristics now exists for this: task
arithmetic~\citep{ilharco2023task}, sign-election methods such as
TIES~\citep{yadav2023ties}, stochastic sparsification such as
DARE~\citep{yu2024dare}, subspace alignment such as
KnOTS~\citep{stoica2025knots}, and closed-form least-squares approaches in
the RegMean line~\citep{jin2023regmean,lorm2024,regmeanpp2025}.

These methods are evaluated almost entirely by benchmark averages, and the
resulting picture is unstable: the same method can be best on one cohort of
adapters and worst on another, with no way to tell in advance which
situation one is in. This paper is about one property of a cohort that
explains a good deal of that instability, and what we contribute is a
\emph{pre-merge diagnostic} for it: the property is measurable from the
adapter weight files alone, before any merge is attempted and in about a
minute of CPU time, and it predicts which family of merging methods will
survive the cohort and how much regularisation a solve-based one needs on
it. It is independent of which merging method is chosen afterwards, which is
what makes it worth running first.

\paragraph{The property.} LoRA factorises each update as $BA$, initialises
$A$ at random and $B$ at zero, and trains. Whether the $T$ adapters of a
cohort drew \emph{one} $A$ or $T$ independent ones is a detail almost never
reported, and it is what a single global training seed silently decides.
When the draw is shared, the tasks end up occupying nearly the same
subspace. Across four base models, sharing the draw moves the median
principal cosine between task row spaces from $0.047$ to $0.995$
(\S\ref{sec:regimes}, Table~\ref{tab:geometry}). The task subspaces stop
being distinguishable directions and become one direction measured four
times. Nothing about this is a bug in LoRA, and nothing about it is visible
in a training curve or a benchmark average.

\paragraph{Measuring it is cheap.} The principal cosines are computed from
the adapter weight files alone: no forward passes, no evaluation data, no
merging, and about one minute of CPU time for a cohort of four rank-16
adapters (\S\ref{subsec:audit}). This is the part of the paper we expect
someone to use directly, and the rest of the paper is an argument that it
is worth running.

\paragraph{The effect splits cleanly by method.} We pre-registered the test
and report both halves (\S\ref{sec:methods}). Across twenty
base-by-method cells spanning five methods, task arithmetic, TIES, DARE,
TVQ and a repaired KnOTS, fifteen show no detectable effect at our $0.01$
nat resolution, five favour independent initialisation and none favours
shared. Scored instead on GSM8K and HumanEval, a pre-registered replication
of the same comparison finds no effect in any of forty cells. The fifteen
have a mechanism: each of these methods returns a bounded-norm near-mean
solution that never enters the ill-conditioned directions, and its energy
there is a rounding error. Of the five cells
that do clear the gate, three are the subspace-alignment method, which is
the heuristic most exposed to a change in subspace geometry.

Methods that solve the linear system behave differently. They reach $35$ to
$57$ times the mean task delta with three quarters of their energy exactly
in the ill-conditioned directions, and run unregularised they are $2.1$ to
$65$ times worse on a shared-initialisation cohort, $25$ to $117$ times
with rank truncation removed. A single Tikhonov term repairs it. We show
this for our own encoder and for RegMean, so it is a property of the
family rather than of our construction.

\textbf{The obvious objection is that regularising a least-squares solve is
standard practice, so we checked what the standard practice actually does.}
RegMean ships a regulariser of its own, an $\alpha$-shrinkage of the
off-diagonal Gram entries with a default of $\alpha = 0.9$. On a collapsed
cohort it is equivalent to a ridge of about $5\times10^{-4}$, roughly $20$
times smaller than the $\lambda^\star = 0.01$ the sweep selects, and it leaves
the system at condition number $6.6$ to $9.1\times10^{3}$ against
$1.5\times10^{2}$ with the ridge. It applies the same dose on a
well-conditioned cohort, where it is ample. That is the practical content of
this paper in one sentence: \textbf{the default regularisation does not know
which regime it is in, and the audit is what tells you}
(\S\ref{subsec:regmean-default}).

\paragraph{The regime is common in released adapters.} Every cohort
discussed so far is one we trained ourselves, in both conditions,
deliberately. So we ran the same measurement over public LoRA cohorts under
a sampling frame registered before any adapter was downloaded: walk the
HuggingFace adapter index by downloads descending, group by publisher and
base model, take the first $50$ cohorts meeting criteria that are all
checkable from metadata, skip none after inspection. \textbf{Seventeen of
the $45$ scorable cohorts, $37.8\%$, are in the collapsed regime}
(\S\ref{subsec:prevalence}). The distribution is bimodal rather than a
continuum: one cohort of $45$ lies between the two populations.

\paragraph{It reaches accuracy, not just our metric.} A large difference
measured in a loss nobody deploys would still be open to the objection that
the loss is the wrong unit. On the same merges, scored on GSM8K and
HumanEval with the repaired scorers, the ridge recovers the shared arm past
the noise gate on four of four bases on both benchmarks. On Mistral the
unregularised shared-cohort merge scores \textbf{zero} on both, with all
$664$ generations non-empty and none discarded, and recovers to $0.46$ and
$0.24$ (\S\ref{subsec:downstream-cond}).

\paragraph{Why conditioning is the right quantity.} The task subspaces
being near-collinear does not by itself make merging impossible; $T$
rank-$r$ adapters with $Tr \leq d$ stay linearly independent even when
nearly collinear, so an exact interpolator always exists. What changes is
its norm. Under a shared draw the exact interpolator is $43$ to $64$ times
larger than the mean task delta; under independent draws it is $T$, the
orthogonal-subspace value, which the measurement returns to within $1.4\%$
at both $T = 3$ and $T = 4$. The condition number of the operator $\bar H$
built from the task subspaces moves correspondingly, from about $1.5$ to
between $1.8\times10^{4}$ and $8.3\times10^{4}$ (\S\ref{sec:regimes},
Figure~\ref{fig:two-regimes}). A method that solves for that interpolator
inherits the conditioning; a method that averages does not. That is the
whole of the split reported above.

We report the principal cosines rather than $\kappa(\bar H)$ throughout,
because the cosines are properties of the row spaces alone and hold under
any curvature model, whereas $\kappa(\bar H)$ depends on the specialisation
$H_t = P_{V_t}$ we adopt (\S\ref{subsec:conditioning}).

\paragraph{Where this came from.} We arrived at the conditioning question
through a rate-distortion analysis of merging, which is in the paper because
it is the reason we looked at $\bar H$ at all and for no stronger reason than
that. \S\ref{sec:setup} says what it did and did not do, and
Appendix~\ref{app:rd} carries the formal development. In short: its floor
term is exactly zero on every cohort we measured, its rate term rests on an
assumption we state but do not prove and on an imported constant, and it
derives none of the four results above. Everything in \S\ref{sec:regimes} and
\S\ref{sec:methods} stands or falls on its own measurements.

\paragraph{What did not survive.} Several pre-registered tests went against
us, and the results below are what is left after them. Our own encoder's
advantage does not survive the move from shared-initialisation cohorts to
properly initialised ones. A control we designed ourselves withdrew our
earlier headline claim that published merging benchmarks are confounded by
initialisation geometry, and it stays withdrawn: a later margin-aware
control would reinstate it, but that control was written after its data had
been seen and we do not use it that way (\S\ref{subsec:withdrawn}). The rate exponent the bound predicts is unresolved:
our sweep lacks the dynamic range to measure it (\S\ref{subsec:w3}). One of our four adapters never learned its
task, so our cohorts are $T=4$ nominal and $3$ effective, and we report
every geometric quantity at both. Appendix~\ref{app:prereg} summarises all
twelve pre-registrations and gives the commit ordering that makes them
checkable;
the documents themselves are in the public repository.

\paragraph{An earlier version of this work is public.} It is posted on Zenodo
under our own names, titled \emph{A Rate-Distortion Function for Model
Merging} \citep{pathak2026rd}, and it predates the pre-registered replication
reported here. We cite it rather than leave it to be found, because it still
presents as headline results two things this paper does not claim: that the
rate-distortion bound is the contribution, and that our encoder attains the
lowest worst-task loss on all four bases at default hyperparameters. Both are
superseded by the tests in \S\ref{sec:methods}. Where the two documents
disagree, this one is current.

\paragraph{Contributions.}
\begin{enumerate}
\item \textbf{A one-CPU-minute audit} that reports, from the adapter weight
  files alone, whether a cohort's task subspaces have collapsed, and hence
  whether an unregularised solve-based merge will fail on it
  (\S\ref{subsec:audit}).
\item \textbf{A pre-registered prevalence estimate} for that regime in
  adapters we did not train: under a sampling frame fixed before any
  download, $37.8\%$ of $45$ public LoRA cohorts are collapsed, and the
  distribution is bimodal rather than continuous
  (\S\ref{subsec:prevalence}).
\item \textbf{A pre-registered demonstration that the collapse is invisible
  to standard merging heuristics and decisive for solve-based ones}, for
  our encoder and for RegMean, at $n = 3$ cohorts in \emph{both} regimes,
  with the mechanism measured: the heuristics never place more than a
  rounding error of their norm in the ill-conditioned directions and the
  unregularised solvers place three quarters of it there
  (\S\ref{subsec:ridge}, \S\ref{subsec:shared-arm}). The invisible half
  is measured twice, in worst-task NLL excess and in downstream accuracy,
  because the metric the first uses is one we ourselves criticise: zero of
  forty accuracy cells show an effect, at a resolution of about six points
  on GSM8K and eleven on HumanEval.
\item \textbf{A pre-registered demonstration that this reaches downstream
  accuracy}, so it is not an artifact of the loss we measure it in: on
  Mistral the unregularised shared-cohort merge scores zero on GSM8K and on
  HumanEval, with no scorer discards, and one ridge term recovers it
  (\S\ref{subsec:downstream-cond}).
\end{enumerate}

The rate-distortion bound (Theorems~\ref{thm:general} and~\ref{thm:achv})
is provenance rather than a contribution, for the reasons above. We also
report a negative result for DARE composed with task arithmetic and the
pre-registered test showing it does not extend to DARE with TIES
(\S\ref{subsec:dare}), which bounds an earlier claim of ours rather than
adding a new one.
